% Options for packages loaded elsewhere
\PassOptionsToPackage{unicode}{hyperref}
\PassOptionsToPackage{hyphens}{url}
\documentclass[
]{article}
\usepackage{xcolor}
\usepackage[margin=1in]{geometry}
\usepackage{amsmath,amssymb}
\setcounter{secnumdepth}{-\maxdimen} % remove section numbering
\usepackage{iftex}
\ifPDFTeX
  \usepackage[T1]{fontenc}
  \usepackage[utf8]{inputenc}
  \usepackage{textcomp} % provide euro and other symbols
\else % if luatex or xetex
  \usepackage{unicode-math} % this also loads fontspec
  \defaultfontfeatures{Scale=MatchLowercase}
  \defaultfontfeatures[\rmfamily]{Ligatures=TeX,Scale=1}
\fi
\usepackage{lmodern}
\ifPDFTeX\else
  % xetex/luatex font selection
\fi
% Use upquote if available, for straight quotes in verbatim environments
\IfFileExists{upquote.sty}{\usepackage{upquote}}{}
\IfFileExists{microtype.sty}{% use microtype if available
  \usepackage[]{microtype}
  \UseMicrotypeSet[protrusion]{basicmath} % disable protrusion for tt fonts
}{}
\makeatletter
\@ifundefined{KOMAClassName}{% if non-KOMA class
  \IfFileExists{parskip.sty}{%
    \usepackage{parskip}
  }{% else
    \setlength{\parindent}{0pt}
    \setlength{\parskip}{6pt plus 2pt minus 1pt}}
}{% if KOMA class
  \KOMAoptions{parskip=half}}
\makeatother
\usepackage{longtable,booktabs,array}
\newcounter{none} % for unnumbered tables
\usepackage{calc} % for calculating minipage widths
% Correct order of tables after \paragraph or \subparagraph
\usepackage{etoolbox}
\makeatletter
\patchcmd\longtable{\par}{\if@noskipsec\mbox{}\fi\par}{}{}
\makeatother
% Allow footnotes in longtable head/foot
\IfFileExists{footnotehyper.sty}{\usepackage{footnotehyper}}{\usepackage{footnote}}
\makesavenoteenv{longtable}
\usepackage{graphicx}
\makeatletter
\newsavebox\pandoc@box
\newcommand*\pandocbounded[1]{% scales image to fit in text height/width
  \sbox\pandoc@box{#1}%
  \Gscale@div\@tempa{\textheight}{\dimexpr\ht\pandoc@box+\dp\pandoc@box\relax}%
  \Gscale@div\@tempb{\linewidth}{\wd\pandoc@box}%
  \ifdim\@tempb\p@<\@tempa\p@\let\@tempa\@tempb\fi% select the smaller of both
  \ifdim\@tempa\p@<\p@\scalebox{\@tempa}{\usebox\pandoc@box}%
  \else\usebox{\pandoc@box}%
  \fi%
}
% Set default figure placement to htbp
\def\fps@figure{htbp}
\makeatother
\setlength{\emergencystretch}{3em} % prevent overfull lines
\providecommand{\tightlist}{%
  \setlength{\itemsep}{0pt}\setlength{\parskip}{0pt}}
\usepackage{bookmark}
\IfFileExists{xurl.sty}{\usepackage{xurl}}{} % add URL line breaks if available
\urlstyle{same}
% fallback for those not using the hyperref driver hyperxmp:
\makeatletter
\@ifundefined{xmpquote}{\newcommand{\xmpquote}[1]{#1}}{}
\makeatother
\hypersetup{
  pdftitle={Exercises},
  hidelinks,
  pdfcreator={LaTeX via pandoc}}

\title{Exercises}
\author{}
\date{\vspace{-2.5em}}

\begin{document}
\maketitle

~

\subsection{Linear models - Model
selection}\label{linear-models---model-selection}

~

In the previous exercise you fitted a pre-conceived model which included
the main effects of the area of the forest patch (\texttt{LOGAREA}), the
grazing intensity (\texttt{FGRAZE}) and the interaction between these
two explanatory variables (\texttt{FGRAZE:LOGAREA}). This was useful as
a training exercise, and might be a viable approach when analysing these
data if an experiment had been designed to test these effects only.
However, if other potentially important variables are not included in
the model this may lead to biased inferences (interpretation).
Additionally, if the goal of the analysis is to explore what models
explain the data in a parsimonious way (as opposed to formally testing
hypotheses), we would also want to include relevant additional
explanatory variables.

~

Here we revisit the previous loyn data analysis, and ask if a `better'
model for these data could be achieved by including additional
explanatory variables and by performing model selection. Because we
would like to test the significance of the interaction between
\texttt{LOGAREA}, and \texttt{FGRAZE}, whilst accounting for the
potential effects of other explanatory variables, we will also include
\texttt{LOGAREA}, \texttt{FGRAZE} and their interaction
(\texttt{FGRAZE:LOGAREA}) in the model as before. Including other
interaction terms between other variables may be reasonable, but we will
focus only on the \texttt{FGRAZE:LOGAREA} interaction as we have
relatively little information in this data set (67 observations). This
will hopefully avoid fitting an overly complex model which will estimate
many parameters for which we have very little data. This is a balance
you will all have to maintain with your own data and analyses (or better
still, perform a power analysis before you even collect your data). No
4-way interaction terms in your models please!

~

It's also important to note that we will assume that all the explanatory
variables were collected by the researchers because \emph{they believed}
them to be biologically relevant for explaining bird abundance
(i.e.~data were collected for a reason). Of course, this is probably not
your area of expertise but it is nevertheless a good idea to pause and
think what might be relevant or not-so relevant and why. This highlights
the importance of knowing your study organism / study area and
discussing research designs with colleagues and other experts in the
field before you collect your data. What you should try to avoid is
collecting heaps of data across many variables (just because you can)
and then expecting your statistical models to make sense of it for you.
As mentioned in the lecture, model selection is a relatively
controversial topic and should not be treated as a purely mechanical
process (chuck everything in and see what comes out). Your expertise
needs to be woven into this process otherwise you may end up with a
model that is implausible or not very useful (and all models need to be
useful!).

~

1. Import the `loyn.txt' data file into RStudio and assign it to a
variable called \texttt{loyn}. Here we will be using all the explanatory
variables to explain the variation in bird density. If needed, remind
yourself of your data exploration you conducted previously. Do any of
the remaining variables need transforming (i.e.~\texttt{AREA},
\texttt{DIST}, \texttt{LDIST}) or converting to a factor type variable
(i.e.~\texttt{GRAZE})? Add the transformed variables to the
\texttt{loyn} dataframe.

~

2. Let's start with a very quick graphical exploration of any potential
relationships between each explanatory variable (collinearity) and also
between our response and explanatory variables (what we're interested
in). Create a pairs plot using the function \texttt{pairs()}of your
variables of interest. Hint: restrict the plot to the variables you
actually need. An effective way of doing this is to store the names of
the variables of interest in a vector
\texttt{VOI\ \textless{}-\ c("Var1",\ "Var2",\ ...)} and then use the
naming method for subsetting the data set \texttt{Mydata{[},\ VOI{]}}.
If you feel like it, you can also add the correlations to the lower
triangle of the plot as you did previously (don't forget to define the
function first).

~

3. Now, let's fit our maximal model. Start with a model of
\texttt{ABUND} and include all explanatory variables as main effects.
Also include the interaction \texttt{LOGAREA:FGRAZE} but no other
interaction terms as justified in the preamble above. Don't forget to
include the transformed versions of the variables where appropriate (but
not the untransformed variables as well otherwise you will have very
strong collinearity between these variables!). Perhaps, call this model
\texttt{M1}.

~

4. Have a look at the summary table of the model using the
\texttt{summary()} function. You'll probably find this summary is quite
complicated with lots of parameter estimates (14) and P values testing
lots of hypotheses. Are all the P values less than our cut-off of 0.05?
If not, then this suggests that some form of model selection is
warranted to simplify our model.

Before you go any further though, a word of caution about looking at 14
tests all at once. If none of these explanatory variables had any real
effect on bird abundance whatsoever, roughly how many of the 14 P values
would you still expect to fall below 0.05 purely by chance?

~

5. Let's perform a first step in model selection using the
\texttt{drop1()} function and use an \emph{F} test based model selection
approach. This will allow us to decide which explanatory variables may
be suitable for removal from the model. Remember to use the
\texttt{test\ =\ "F"} argument to perform \emph{F} tests when using
\texttt{drop1()}. Which explanatory variable is the best candidate for
removal and why?

What hypothesis is being tested when we do this model selection step?

~

\textbf{A note on keeping track of model selection.} Describing what you
did during model selection is genuinely difficult to do well in writing,
and it is done badly in the published literature more often than almost
anything else on this course. Papers routinely say ``a minimum adequate
model was obtained by backward selection'' and leave the reader with no
idea what was in the starting model, what came out, in what order, or on
what grounds.

A summary table solves this, and it is worth seeing one being built
rather than just being shown a finished one. So as we work through the
deletions below I'll add a row to a table at each step, and you'll find
it updated in the solution to each question. By the time we reach the
end we'll have something that could go straight into a results section.
You don't need to construct it yourself here, just follow along, but do
consider making one like it when you report a model selection of your
own.

Here is the structure, with the starting model in it:

\textbf{Table 1: Summary of model selection.}

{\def\LTcaptype{none} % do not increment counter
\begin{longtable}[]{@{}
  >{\raggedright\arraybackslash}p{(\linewidth - 10\tabcolsep) * \real{0.0732}}
  >{\raggedright\arraybackslash}p{(\linewidth - 10\tabcolsep) * \real{0.6220}}
  >{\raggedright\arraybackslash}p{(\linewidth - 10\tabcolsep) * \real{0.1585}}
  >{\centering\arraybackslash}p{(\linewidth - 10\tabcolsep) * \real{0.0488}}
  >{\centering\arraybackslash}p{(\linewidth - 10\tabcolsep) * \real{0.0488}}
  >{\centering\arraybackslash}p{(\linewidth - 10\tabcolsep) * \real{0.0488}}@{}}
\toprule\noalign{}
\begin{minipage}[b]{\linewidth}\raggedright
Model
\end{minipage} & \begin{minipage}[b]{\linewidth}\raggedright
Terms in the model
\end{minipage} & \begin{minipage}[b]{\linewidth}\raggedright
Term removed
\end{minipage} & \begin{minipage}[b]{\linewidth}\centering
\emph{df}
\end{minipage} & \begin{minipage}[b]{\linewidth}\centering
\emph{F}
\end{minipage} & \begin{minipage}[b]{\linewidth}\centering
\emph{P}
\end{minipage} \\
\midrule\noalign{}
\endhead
\bottomrule\noalign{}
\endlastfoot
M1 & LOGLDIST + LOGDIST + YR.ISOL + ALT + LOGAREA + FGRAZE +
LOGAREA:FGRAZE & -- & -- & -- & -- \\
\end{longtable}
}

Two things about that first row. The starting model is in the table
because a reader cannot judge a selection procedure without knowing what
it started from: dropping four terms from a model of six is a very
different exercise from dropping four from a model of twenty. And the
\emph{F} and \emph{P} columns hold dashes because nothing was removed in
order to arrive at M1. Those columns describe the deletion that got you
from the previous model to this one, and for the first row there isn't
one.

~

~

6. Update and refit your model and remove the least significant
explanatory variable (from above). Repeat single term deletions with
\texttt{drop1()} again using this updated model. You can update the
model by just fitting a new model without the appropriate explanatory
variable and assign it to a new name (\texttt{M2}). Alternatively you
can use the \texttt{update()} function instead (I show you how to do
this in the solutions).

~

7. Again, update the model to remove the least significant explanatory
variable (from above) and repeat single term deletions with
\texttt{drop1()}.

~

8. Once again, update the model to remove the least significant
explanatory variable (from above) and repeat single term deletions with
\texttt{drop1()}.

~

9. And finally, update the model to remove the least significant
explanatory variable (from above) and repeat single term deletions with
\texttt{drop1()}.

~

10. If all goes well, your final model should be
\texttt{lm(ABUND\ \textasciitilde{}\ LOGAREA\ +\ FGRAZE\ +\ LOGAREA:FGRAZE)}
which you encountered in the previous exercise. Also, you may have
noticed that the output from the \texttt{drop1()} function does not
include the main effects of \texttt{LOGAREA} or \texttt{FGRAZE}. Can you
think why this might be the case?

~

11. Now that you have your final model, you should go through your model
validation and model interpretation as usual. As we have already
completed this in the previous exercise I'll leave it up to you to
decide whether you include it here (you should be able to just copy and
paste the code).

Please make sure you understand the biological interpretation of each of
the parameter estimates and the interpretation of the hypotheses you are
testing. As before, lead with the estimate and its confidence interval
(\texttt{confint()}) and treat the P value as supporting evidence rather
than as the finding
(\href{https://alexd106.github.io/EV5020/linear_model_1_exercise.html\#pval_note}{reminder}).

There is also one caveat that applies to your final model here but did
not apply to the model you fitted in the previous exercise, even though
the two models are identical. Have a think about it before you look at
the solutions. The P values and confidence intervals for your final
model are calculated as though you had specified that model in advance.
You didn't: you arrived at it by looking at these same data. Does that
make them too optimistic, too conservative, or does it make no
difference at all?

~

\textbf{One last caution, and it applies to everything you have done in
this exercise.} These are observational data. Nobody assigned patch
areas or grazing intensities to these forests, they were measured as
they were found. A linear model will happily estimate an effect of
grazing on bird abundance, but it cannot tell you that grazing
\emph{causes} that change.

Two things are worth keeping in mind. First, variables you didn't
measure may be doing the work. Heavily grazed patches might also be the
ones nearest roads, or on more fertile soil, or with a different
understorey, and any of those could be the real driver of the difference
in bird abundance.

Second, your explanatory variables are correlated with one another, as
the pairs plot in Question 2 showed. That has a specific consequence for
how you read the output: every estimate in your model is the effect of
that variable \emph{holding the other variables in the model constant}.
Change which variables are in the model and the estimate changes, as you
saw during model selection. So a coefficient is never simply ``the
effect of grazing''; it is the effect of grazing after accounting for
whatever else you happened to include, and it will absorb some of the
effect of any correlated variable you left out.

None of this means observational studies aren't worth doing, and most of
ecology is observational for perfectly good reasons. It does mean your
language should match the evidence you actually have. ``Bird abundance
was lower in more heavily grazed patches'' is supported by these data.
``Grazing reduces bird abundance'' is a stronger claim than a model like
this one can carry on its own, and it needs experimental evidence, or a
well understood mechanism, standing behind it.

~

12. Now for reporting your results, and this is the last write-up of the
course. You arrived at your final model by a selection procedure, which
brings a couple of reporting problems the previous three exercises
didn't have.

This is where the table you have been building since Question 5 earns
its keep. Present it as a table in your results, exactly as you would in
a paper, and then \textbf{refer to it} in your text rather than
describing every deletion in prose. That is what the table is for: it
carries the detail, so your paragraph can carry the argument.

Write \textbf{no more than 200 words} of accompanying text for the
methods and results of a paper. Your reader is an ecologist who knows
their stuff but isn't a statistician, and who wasn't in the room while
you did the analysis. They need to be able to work out from your text
and your table exactly what you did and how much to trust it. Questions
4 and 11 have already told you the two things an honest account has to
include that a naive one would leave out.

~

\textbf{OPTIONAL questions} if you have time / energy / inclination!

~

A1. If we weren't aiming to directly test the effect of the
\texttt{LOGAREA:FGRAZE} interaction statistically (i.e.~test this
specific hypothesis), we could use AIC to perform model selection.
Repeat the model selection you did above, but this time use the
\texttt{drop1()} function and perform model selection using AIC instead.
Don't forget, if we want to perform model selection based on AIC with
the \texttt{drop1()} function we need to omit the \texttt{test\ =\ "F"}
argument)

~

\textbf{A table again, and two ways of reading AIC.} As with the
\emph{F} test route, I'll build a summary table as we go. Table 2 has
the same shape as Table 1, but with AIC in place of the \emph{F}
statistic and \emph{P} value, and with two AIC columns rather than one.
Those two columns carry the same evidence read in opposite directions,
so it is worth being clear about them before you see any numbers.

\textbf{AIC change} is what happened to AIC when a term was deleted.
This is the number the selection decision is actually based on, and it
is what \texttt{drop1()} puts in front of you. A negative value means
AIC fell when the term came out, so the simpler model was preferred.

\textbf{ΔAIC} is the column you will see in published tables: each
model's AIC minus the \textbf{lowest} AIC in the set. The preferred
model therefore has ΔAIC = 0, and every other model has a positive
value.

The usual rule of thumb (\hyperref[refs]{Burnham and Anderson 2002})
applies to both columns, but reads differently in each:

\begin{itemize}
\tightlist
\item
  \textbf{AIC change:} a change of less than 2 either way is no real
  support for preferring either model, so the simpler one wins and the
  term goes. A deletion that \emph{raises} AIC by roughly 4 to 7 is
  considerable support for keeping the term, and by more than 10, strong
  support.
\item
  \textbf{ΔAIC:} a model within 2 of the best has substantial support,
  one 4 to 7 above the best has considerably less, and one more than 10
  above has essentially none.
\end{itemize}

\emph{k}, in the third column, is simply the number of parameters the
model estimates.

~

~

A2. Refit your model with the variable associated with the lowest AIC
value removed. Run \texttt{drop1()} again on your updated model. Perhaps
call this new model \texttt{M2.AIC}.

~

A3. Refit your model with the variable associated with the lowest AIC
value removed and run \texttt{drop1()} again on your new model
(\texttt{M3.AIC}).

~

A4. Repeat your model selection by removing the variable indicated by
the model with the lowest AIC.

~

A5. Rinse and repeat as above.

~

If all goes well, your final model should be
\texttt{lm(ABUND\ \textasciitilde{}\ LOGAREA\ +\ FGRAZE\ +\ LOGAREA:FGRAZE)}.
This is the same model you ended up with when using the \emph{F} test
based model selection. This might not always be the case and generally
speaking AIC based model selection approaches tend to favour more
complicated minimum adequate models compared to \emph{F} test based
approaches.

We don't need to re-validate or re-interpret the model, since we have
already done this previously.

You have already seen how to present this. Table 2 above is the AIC
equivalent of Table 1, and either one can go straight into a results
section with the text referring to it, exactly as in Question 12.

A couple of practical notes if you build one of these for your own work.
The AIC values reported by \texttt{drop1()} are not the same as those
from the \texttt{AIC()} function, because the two calculate AIC on
slightly different scales. Neither is wrong, but never mix values from
the two in the same table. And if you want a formatted table in a report
rather than console output, \texttt{knitr::kable()} will do it (see the
\href{https://intro2r.com/install-rm.html}{`Installing R Markdown'}
appendix of the Introduction to R book).

~

\subsubsection{References}\label{refs}

Burnham, K.P. and Anderson, D.R. (2002) \emph{Model Selection and
Multimodel Inference: A Practical Information-Theoretic Approach}, 2nd
edition. Springer, New York.

~

End of the model selection exercise

\end{document}
